% !TEX program = xelatex

\documentclass{resume}
%\usepackage{zh_CN-Adobefonts_external}

\begin{document}
\pagenumbering{gobble}

\name{张钧瑞 \\ Junrui Zhang}

\basicInfo{
  \email{your.email@example.com} \textperiodcentered\ 
  \phone{(+86) xxx-xxxx-xxxx} \textperiodcentered\ 
  \github[Zjrua]{https://github.com/Zjrua}}

\section{\faGraduationCap\ 教育背景}
\datedsubsection{\textbf{哈尔滨工业大学 (HIT)}, 哈尔滨}{2025 -- 2028}
\textit{硕士研究生}，应用统计 (Applied Statistics)
\datedsubsection{\textbf{哈尔滨工业大学 (HIT)}, 哈尔滨}{2020 -- 2024}
\textit{理学学士}，信息与计算科学 (Information and Computing Science)

\section{\faFlask\ 科研经历}
\datedsubsection{\textbf{体质数据分析研究}}{2025.09 -- 至今}
\role{Python, PyTorch, 统计分析}{导师课题组}
对体质健康数据进行统计建模与分析，论文已投稿《体育科学》。
\begin{itemize}
  \item 主导数据处理与统计建模全流程
  \item 应用线性回归、多元分析等方法探索体质影响因素
\end{itemize}

\section{\faTrophy\ 竞赛获奖}
\datedsubsection{\textbf{全国统计建模大赛}（研究生组）}{2026}
\role{参赛编号: TJJM20260414180979}{国家级}
\datedline{\textit{美国大学生数学建模竞赛 (MCM/ICM) Meritorious Winner}}{2022}
\datedline{\textit{全国大学生数学建模竞赛 黑龙江省二等奖}}{2023}

\section{\faCogs\ 技能}
\begin{itemize}[parsep=0.5ex]
  \item 编程语言：Python > C/C++ > SQL
  \item 框架与库：PyTorch, pandas, NumPy, scikit-learn
  \item 工具：Git, Linux, Jupyter, VS Code, \LaTeX
  \item 语言：中文（母语），英语（流利）
\end{itemize}

\section{\faInfo\ 其他}
\begin{itemize}[parsep=0.5ex]
  \item GitHub: \href{https://github.com/Zjrua}{https://github.com/Zjrua}
  \item 高中学过信息学竞赛 (NOI)
  \item 入党积极分子
\end{itemize}

\end{document}
